\documentclass[11pt]{amsart}
\usepackage{geometry}                % See geometry.pdf to learn the layout options. There are lots.
\geometry{letterpaper}                   % ... or a4paper or a5paper or ... 
%\geometry{landscape}                % Activate for for rotated page geometry
%\usepackage[parfill]{parskip}    % Activate to begin paragraphs with an empty line rather than an indent
\usepackage{graphicx}
\usepackage{amssymb}
\usepackage{epstopdf}
\DeclareGraphicsRule{.tif}{png}{.png}{`convert #1 `dirname #1`/`basename #1 .tif`.png}

\title{PROBLEMA 3.17}
%\date{}                                           % Activate to display a given date or no date

\begin{document}
\maketitle
%\section{}
%\subsection{}
\noindent Se dice que un entero positivo N es un número primo si los únicos enteros positivos que lo dividen son exactamente 1 y N. Diseñe un diagrama de flujo que lea un número M, y obtenga y cuente todos los números primos menores a M.\\

\noindent  Dato: M (variable de tipo entero que representa el número límite que se ingresa).\\

\noindent \textbf {Explicación de las variables} \\

\noindent \quad I: Variable de tipo entero. Representa la variable de control del ciclo externo. Se inicializa en 3 porque sabemos por definición que el número 1 y los números pares (con excepción del 2) no son primos, y se incrementa de 2 en 2. \\

\noindent \quad SP: Variable de tipo entero. Acumula el número de números primos comprendidos entre 1 y M.\\

\noindent \quad M: Variable de tipo real.\\

\noindent \quad  J: Variable de tipo entero. Representa la variable de control del ciclo interno. Se inicializa en 3 y se incrementa de 2 en 2.\\

\noindent \quad BAND: Variable de tipo caracter. Se inicializa en verdadero (‘V’). Si alguna de las divisiones que se realizan produce módulo cero, entonces toma el valor de falso (‘F’).\\
\end{document}  
